% Part: computability
% Chapter: computability-theory
% Section: ce-sets

\documentclass[../../../include/open-logic-section]{subfiles}

\begin{document}

\olfileid{cmp}{thy}{ces}
\olsection{قابلِ حساب گنتی وار فہرست بنن جوگ سیٹ}

\begin{defn}
کوئی سیٹ \emph{قابلِ حساب گنتی وار فہرست بنن جوگ
(computably enumerable)} اُدوں اے جدوں اوہ خالی
ہووے یا کسے قابلِ حساب فنکشن دیاں حاصل قدراں دا
سیٹ ہووے۔
\end{defn}

\begin{history}
قابلِ حساب گنتی وار فہرست بنن جوگ سیٹاں نوں
\emph{بازگشتی طور اُتے گنتی وار فہرست بنن جوگ
(recursively enumerable)} سیٹ وی آکھیا جاندا اے۔
ایہ اصلی اصطلاح اے؛ اج کل دوہاں اصطلاحاں دے نال
نال اوہناں دے مخفف ``c.e.'' تے ``r.e.'' وی عام
ورتے جاندے نیں۔
\end{history}

\begin{explain}
تعریف دا مطلب کی اے تے ایہ اصطلاح کیوں مناسب اے،
ایس اُتے غور کرو۔ خیال ایہ اے کہ جے $S$ قابلِ حساب
فنکشن~$f$ دیاں حاصل قدراں دا سیٹ ہووے، تاں
\[
S = \{ f(0), f(1), f(2), \dots \},
\]
ایس لئی $f$ نوں $S$ دے رکناں دی گنتی وار فہرست
بناؤن والا فنکشن سمجھیا جا سکدا اے۔ دھیان دیو،
تعریف مطابق $f$ دا ودھدا جان والا فنکشن ہونا ضروری
نہیں؛ یعنی فہرست ودھدی ترتیب وچ ہونی لازمی نہیں۔
بلکہ $f$ دا اک بہ اک (injective) ہونا وی ضروری نہیں؛
فہرست $f(0)$، $f(1)$، $f(2)$، \dots{} وچ $S$ دے
رکناں دی دہرائی ہو سکدی اے۔ مثال لئی مستقل فنکشن
$f(x) = 0$ سیٹ $\{ 0
\}$ دی گنتی وار فہرست بناندا اے۔
\end{explain}

ہر قابلِ حساب سیٹ قابلِ حساب گنتی وار فہرست بنن جوگ
ہوندا اے۔ ایہ ویکھن لئی من لو $S$ قابلِ حساب اے۔
جے $S$ خالی اے، تاں تعریف مطابق اوہ قابلِ حساب
گنتی وار فہرست بنن جوگ اے۔ نہیں تاں $a$ نوں $S$
دا کوئی وی رکن چن لو۔ $f$ دی تعریف ایویں کرو:
\[
f(x) =
\begin{cases}
x & \text{جے $\Char{S}(x) = 1$} \\
a & \text{ہور صورت وچ۔}
\end{cases}
\]
تاں $f$ قابلِ حساب فنکشن اے، تے $S$ فنکشن~$f$
دیاں حاصل قدراں دا سیٹ اے۔

\end{document}
